\begin{tabular}{p{0.44\textwidth} p{0.48\textwidth}}
\toprule
\textbf{What is proved} & \textbf{What the simulation adds} \\
\midrule
The weighted sum returns one verdict on any case whatever. A rule certifying each direction separately returns two whenever at least two principles apply strongly each way, which is a condition on the case and not on the rule. & on the directly opposed battery, where that condition holds by construction, the rule contradicts itself on 100\% of cases and the sum on 0\%; the condition itself obtains on 45\% of independent draws, 11\% of category-correlated ones, and 0\% under full correlation \\
\addlinespace
In the reliance mechanism, conferred weight is a function of pressure and activations and not of invocation propensity. Not proved for uptake, where weight is invocation frequency by construction. & two runs differing only in invocation frequency confer identical weights, to a difference of 0.000, though their uptake shares differ \\
\addlinespace
The weighted sum responds to a change in magnitude alone; a rule reading only ranks cannot. & magnitude sensitivity 0.023 for the sum against exactly 0.000 for a rank-only rule with no activation cutoff, and 0.0024 for the implemented foil, which carries one \\
\addlinespace
For activations drawn independently and uniformly, net force has variance $\tfrac{1}{3}\sum_i w_i^2$ and the determinacy rate is $2(1-\Phi(\tau/\sigma))$ under a normal approximation. Not proved for correlated activations. & simulation matches the formula to within 0.006 on the grid tested here, and 0.0106 on the finer grid of Part VIII; both residuals measure the approximation rather than the implementation \\
\addlinespace
With equal weights the sum and a vote-count disagree on a set of positive measure. The proof is a counterexample and supplies no rate. & on nine principles the two disagree on 15\% of cases, reproduced at 15.3\% on the separate implementation of Part VII \\
\bottomrule
\end{tabular}
